\chapter{소프트웨어 구현}\label{ch:software}

소프트웨어 모델은 이 프로젝트의 \emph{기준 모델}이다. 기계식, 릴레이식, 전자식(FPGA)의 세 기계는 모두
이 모델이 내놓는 결과와 일치해야 한다. 드럼 한 벌이나 릴레이 접점망이 올바른 치환을 내는지, 각 기계가
올바른 위치에서 정지하는지를 판정할 근거가 이 모델이다. 그래서 모델은 답만 맞추는 탐색기가 아니라 \ref{ch:bombe}장의 기계를
구성요소 단위로 옮긴 것이어야 한다. 이 장은 그 구조와 대응 관계, 그리고 구현하면서 얻은 교훈을 적는다.

\section{전체 구조}

코어는 Rust로 쓴 라이브러리(\code{christopher-core})이고, 사람이 다루는 인터페이스는 Python
패키지(\code{christopher})다. 둘 사이에는 얇은 바인딩 계층이 있다(그림~\ref{fig:arch}).

\begin{figure}[ht]
\centering
\begin{tikzpicture}[x=1cm,y=1cm,node distance=6mm]
  \node[blkg,minimum width=58mm] (py) at (0,3.2) {\code{python/christopher}\\{\footnotesize 조작 API: \code{Menu}, \code{Bombe}, \code{Stop}, \code{Solution}}};
  \node[blkn,minimum width=58mm] (bind) at (0,1.9) {\code{christopher-py} (PyO3 바인딩)};
  \node[blkb,minimum width=58mm,minimum height=16mm] (core) at (0,0.2)
     {\code{christopher-core} (순수 Rust)\\[1pt]{\footnotesize\code{enigma} · \code{menu} · \code{bombe} · \code{check} · \code{search}}};
  \node[blko,minimum width=28mm] (wir) at (-6.2,0.2) {\code{data/wirings}\\배선표 (원본 하나)};
  \node[blko,minimum width=28mm] (vec) at (6.2,0.2) {\code{data/vectors}\\검증 벡터};
  \node[blkn,minimum width=28mm] (rtl) at (6.2,-1.6) {\code{rtl/}\\디지털 회로};
  \node[blkn,minimum width=28mm] (hw) at (6.2,-3.0) {\code{hardware/}\\실물};
  \node[blk,minimum width=22mm] (user) at (-6.2,3.2) {조작자};
  \draw[->] (user) -- (py);
  \draw[->] (py) -- (bind);
  \draw[->] (bind) -- (core);
  \draw[->] (wir) -- (core);
  \draw[->] (core) -- node[above,font=\scriptsize,text=cgray] {내보내기} (vec);
  \draw[->] (vec) -- node[right,note] {대조} (rtl);
  \draw[->] (rtl) -- (hw);
  \draw[->,dashed] (wir) |- node[pos=0.75,below,note] {같은 배선표에서 생성} (rtl);
\end{tikzpicture}
\caption{소프트웨어의 계층과 다른 단계와의 관계. 코어는 바인딩에 대해 아무것도 모른다. 회로와 실물은
코어가 내보낸 검증 벡터와 대조된다.}
\label{fig:arch}
\end{figure}

설계 원칙은 세 가지다. 왜 그렇게 정했는지와 함께 적는다.

\begin{itemize}
\item \textbf{코어는 바인딩과 분리한다.} 같은 코어를 Python에서도, 회로 검증의 기준 모델로도 그대로 써야
하기 때문이다. 코어에 Python 전용 코드가 섞이면 ``회로가 일치해야 하는 대상''이 무엇인지 흐려진다.
\item \textbf{Rust 코어가 기준이다.} 회로와 실물은 코어의 출력과 비교해 검증한다. 기준이 둘이면 둘이
어긋날 때 어느 쪽이 맞는지 알 수 없다.
\item \textbf{배선표는 코드가 아니라 데이터다.} 로터와 반사판의 배선은 \code{data/wirings/enigma.txt}
한 곳에만 있고, Rust 코어와 회로 쪽 생성 스크립트가 같은 파일을 읽는다. 손으로 옮겨 적은 사본이 둘이면
언젠가 한 글자가 달라진다.
\end{itemize}

\section{실제 기계와의 대응}\label{sec:mapping}

표~\ref{tab:mapping}은 실제 봄베의 구성요소가 코드의 어디에 있는지를 보인다. 이 대응이 이후 회로와
실물 단계의 설계도가 된다.

\begin{table}[ht]
\centering
\caption{실제 봄베의 구성요소와 소프트웨어 모델의 대응.}
\label{tab:mapping}
\small
\begin{tabularx}{\linewidth}{>{\raggedright\arraybackslash}p{38mm}>{\raggedright\arraybackslash}p{50mm}>{\raggedright\arraybackslash}X}
\toprule
실제 기계 & 코드 & 비고 \\
\midrule
드럼 세 개 한 벌(스크램블러) & \code{WheelOrder::scramblers} & 플러그보드 없는 에니그마의 치환표 \\
26선 케이블 26개 & \code{[u32; 26]} & 케이블마다 26비트, 비트 하나가 선 한 가닥 \\
메뉴 배선(케이블 꽂기) & \code{Plugged::new} & 글자별 이웃 목록 \\
대각선 보드 & \code{propagate}의 대칭 전파 & 선 $(X,y)$와 $(Y,x)$ \\
전류의 확산 & \code{Plugged::propagate} & 변화가 없을 때까지 반복 \\
테스트 레지스터 & \code{live[test\_letter]} & 26비트 \\
드럼 구동과 자리올림 & \code{run\_slice}의 네 겹 반복 & 탐색 순서는 실제와 다르다(본문) \\
정지 & \code{Stop} & 로터 순서, 코어 위치, 레지스터 \\
체킹 머신 & \code{check::stecker\_hypotheses} & 가설 하나에 전류를 흘려 스테커를 읽는다 \\
Hut~6의 수작업 & \code{check::check\_stop}의 나머지 & 링 설정, 남은 스테커 \\
봄베 여러 대의 동시 운용 & \code{search} & 로터 순서를 스레드에 분배 \\
\bottomrule
\end{tabularx}
\end{table}

실제 기계와 다르게 한 곳이 둘 있다. 둘 다 결과(정지의 집합)에는 영향이 없다.

\begin{itemize}
\item \textbf{탐색 순서.} 실제 봄베는 에니그마의 왼쪽 로터에 해당하는 드럼을 가장 빨리 돌린다%
(표~\ref{tab:drums}). 모델은 오른쪽 로터를 가장 안쪽 반복으로 돈다. 이렇게 하면 가운데와 왼쪽 로터가
고정된 동안 스크램블러 표를 한 번만 계산하면 되기 때문이다(\ref{sec:sw-tables}절).
\item \textbf{위치의 표기.} 모델의 정지 위치는 코어 위치를 \texttt{A}=0으로 적은 것이다. 실제 봄베의
지시 드럼은 눈금이 역순이고 드럼 종류별로 글자가 어긋나 있어 다른 글자가 표시된다%
~\cite{sale-ringstellung,ekhall-hallenberg-tutorial}. 역사적 표기로의 변환은 아직 구현하지 않았다.
\end{itemize}

\section{에니그마 모델}

에니그마 모델(\code{enigma})은 봄베의 공격 대상이자 검증 도구다. 로터 한 단의 통과는
식~\eqref{eq:through}를 그대로 옮긴 함수 하나이고, 스크램블러는 이 함수를 오른쪽, 가운데, 왼쪽,
(그리스 로터,) 반사판, 다시 역순으로 이어 붙인 것이다. 키를 누르면 먼저 전진하고 그다음 암호화한다.

전진 규칙은 세 줄이다. 가운데 로터가 노치 위치에 있으면 가운데와 왼쪽이 함께 움직이고, 그렇지 않고
오른쪽 로터가 노치 위치에 있으면 가운데가 움직이며, 오른쪽은 항상 움직인다. 첫 번째 경우가 이중
스테핑이다(그림~\ref{fig:stepping}).

모델의 정확성은 네 가지로 확인했다. 잘 알려진 시험값(로터 I-II-III, 반사판 B, 설정 \texttt{AAA}에서
\texttt{AAAAA}가 \texttt{BDZGO}로 암호화됨), 이중 스테핑과 노치가 둘인 로터의 전진, M4에서 그리스 로터를
\texttt{A}에 둔 얇은 반사판이 3로터 기종의 반사판과 같다는 등가성, 그리고 실제 전문 두 건의
복호화다(\ref{sec:historical}절).

\section{메뉴 생성}

메뉴 모듈(\code{menu})은 조작자가 봄베를 돌리기 전에 종이 위에서 하던 일을 한다.

\paragraph{정렬.} \code{possible\_offsets}는 크립을 암호문의 모든 위치에 놓아 보고, 같은 글자가 겹치는
위치를 뺀다(\ref{sec:crib}절).

\paragraph{메뉴 구성.} \code{Menu::from\_crib}은 크립의 글자마다 선 하나(평문 글자, 암호문 글자, 크립 안의
위치)를 만든다. 테스트 글자는 가장 큰 덩어리 안에서 선이 가장 많이 모이는 글자로 고른다.

\paragraph{루프 수.} 그래프의 독립적인 루프 수는 $(\text{선 수}) - (\text{글자 수}) + (\text{덩어리 수})$다.
덩어리마다 글자들을 잇는 데 필요한 최소한의 선은 (글자 수 $-$ 1)개이고, 그보다 많은 선은 하나마다 루프를
하나씩 닫기 때문이다. 덩어리는 서로소 집합(union-find)으로 센다.

\paragraph{부분 메뉴 선택.} 기계의 스크램블러 수는 한정되어 있으므로 \code{Menu::select}가 선의 수를
제한한 부분 메뉴를 고른다. 방법은 손으로 메뉴를 만들 때와 같은 탐욕법이다. 선이 가장 많은 글자에서
시작해, 이미 고른 글자 두 개를 잇는 선(루프를 닫는 선)이 있으면 먼저 고르고, 없으면 연결이 많은 새
글자로 뻗는다. 루프를 우선하는 이유는 루프 하나가 거짓 정지를 약 $1/26$로 줄이기 때문이다.
예제 전문에서 크립을 \texttt{WETTERVORHERSAGEFUERDIE}(23글자)로 늘리면 같은 14개의 선으로 글자 11개,
루프 4개인 메뉴가 골라지고, 정지 수는 157회에서 2회로 줄어든다(\ref{sec:results-search}절).

\section{봄베 모델}

\subsection{상태: 676비트}

기계의 상태는 26개 케이블 × 26선의 통전 여부다. 모델은 이를 32비트 정수 26개(\code{[u32; 26]})로
둔다. $c$번째 정수의 $w$번째 비트가 1이면 케이블 $c$의 $w$번 선이 통전된 것이다(그림~\ref{fig:state}).
이 표현을 고른 이유는 둘이다. 테스트 레지스터가 26선 모두 통전되었는지를 정수 비교 한 번으로 알 수
있고, 뒤에 올 디지털 회로의 676비트 레지스터와 비트 단위로 그대로 대응한다.

\begin{figure}[ht]
\centering
\begin{tikzpicture}[x=2.7mm,y=-2.7mm]
  \def\letters{A,B,C,D,E,F,G,H,I,J,K,L,M,N,O,P,Q,R,S,T,U,V,W,X,Y,Z}
  \fill[clight] (-0.5,-0.5) rectangle (25.5,25.5);
  % 테스트 레지스터(케이블 E) 강조
  \fill[corangebg] (-0.5,3.5) rectangle (25.5,4.5);
  \foreach \r/\c in {0/16,4/21,6/10,7/13,8/20,9/15,11/11,12/3,17/2,19/19,20/8,21/4,22/5,25/1,
                     16/0,10/6,13/7,15/9,3/12,2/17,5/22,1/25} {
    \fill[cblue] (\c-0.5,\r-0.5) rectangle (\c+0.5,\r+0.5);
  }
  \foreach \i in {0,...,26} {
    \draw[white,line width=0.4pt] (\i-0.5,-0.5) -- (\i-0.5,25.5);
    \draw[white,line width=0.4pt] (-0.5,\i-0.5) -- (25.5,\i-0.5);
  }
  \foreach \l [count=\i from 0] in \letters {
    \node[font=\ttfamily\tiny] at (-1.3,\i) {\l};
  }
  \foreach \l [count=\i from 0] in {a,b,c,d,e,f,g,h,i,j,k,l,m,n,o,p,q,r,s,t,u,v,w,x,y,z} {
    \node[font=\ttfamily\tiny] at (\i,-1.2) {\l};
  }
  \node[note,rotate=90] at (-3.3,12.5) {케이블 (글자) = 배열의 원소};
  \node[note] at (12.5,-2.8) {선 (스테커 짝 가설) = 원소 안의 비트};
  \draw[corange,line width=0.8pt] (-0.5,3.5) rectangle (25.5,4.5);
  \node[note,anchor=west,text=corange] at (26.3,4) {\code{live[4]}: 테스트 레지스터};
  \node[note,anchor=west,text width=48mm,align=left] at (26.3,12)
    {파란 칸: 통전된 선.\\[3pt]맞는 드럼 위치에서 맞는 가설 $\Stk(\mathtt{E})=\mathtt{V}$에 전압을 건
     상태다. 메뉴 글자의 케이블마다 선이 하나씩 살아 있고, 대각선 보드가 그 대칭 짝을 살린다. 살아 있는
     선을 읽으면 플러그보드의 일부가 그대로 나온다.};
\end{tikzpicture}
\caption{상태의 표현. 가로줄 하나가 케이블 하나이자 \code{u32} 하나다. 그림의 상태는 예제 키의 스테커
(\texttt{AQ BZ CR DM EV FW GK HN IU JP})와 그림~\ref{fig:menu}의 메뉴에서 정해진다.}
\label{fig:state}
\end{figure}

\subsection{위치의 정의}

봄베의 위치(\code{Position})는 \emph{크립의 첫 글자를 암호화하는 순간}의 코어 위치로 정의한다. 에니그마는
암호화 전에 전진하므로, 전문의 첫 글자에 크립을 놓았다면 이 위치는 시작 위치에서 오른쪽 로터가 한 칸
나아간 값이다. 크립 위치 $k$(0부터)의 선은 오른쪽 코어 위치 $r+k$의 스크램블러를 쓴다. 이 한 칸 차이는
틀리기 쉬워서, 알려진 키로 암호화한 전문에서 정지 위치가 정확히 예상한 값인지 확인하는 시험으로
고정했다.

\subsection{스크램블러 표: 바뀌는 것만 다시 계산한다}\label{sec:sw-tables}

한 위치에서 메뉴의 선마다 26글자짜리 치환표가 필요하다. 이를 위치마다 처음부터 계산하면 선 14개에 대해
매번 로터 여섯 번과 반사판을 지나야 한다. 그런데 메뉴의 모든 스크램블러는 가운데와 왼쪽 드럼의 위치가
같고 오른쪽 드럼만 크립 위치만큼 어긋나 있다. 그러므로 가운데와 왼쪽이 정해지면 필요한 치환표는
오른쪽 드럼의 26가지 위치에 대한 26개가 전부이고, 선은 그중 $(r+k) \bmod 26$번째를 골라 쓰기만 하면 된다
(그림~\ref{fig:tables}).

\begin{figure}[ht]
\centering
\begin{tikzpicture}[x=1cm,y=1cm]
  \draw[rounded corners=2pt,draw=cgray,fill=clight] (0,0) rectangle (13.6,5.4);
  \node[anchor=north west,font=\small] at (0.15,5.35) {반복: 그리스 로터 $g$ (M4만) $\to$ 왼쪽 $l$ $\to$ 가운데 $m$ \quad\textcolor{cgray}{($26^2$ 또는 $26^3$회)}};
  % inner
  \node[blkb,minimum width=40mm,minimum height=12mm] (inner) at (3.2,3.6) {안쪽 치환 \code{inner}\\{\footnotesize 가운데·왼쪽·반사판·왼쪽·가운데}};
  \node[blkb,minimum width=44mm,minimum height=12mm] (tab) at (9.4,3.6) {스크램블러 표 26개\\{\footnotesize 오른쪽 드럼 위치 $0\ldots25$}};
  \draw[->] (inner) -- (tab);
  % innermost loop
  \draw[rounded corners=2pt,draw=cblue,fill=white] (0.4,0.3) rectangle (13.2,2.5);
  \node[anchor=north east,font=\small] at (13.1,2.45) {반복: 오른쪽 $r$ \quad\textcolor{cgray}{(26회)}};
  \node[blkg,minimum width=36mm] (pick) at (2.7,1.1) {선마다 표 고르기\\$(r+k) \bmod 26$};
  \node[blkg,minimum width=30mm] (prop) at (7.0,1.1) {전류 확산\\\code{propagate}};
  \node[blkr,minimum width=30mm] (stop) at (11.1,1.1) {레지스터 $\neq$ 26선\\$\Rightarrow$ 정지 기록};
  \draw[->] (pick) -- (prop);
  \draw[->] (prop) -- (stop);
  \draw[->] (tab.south) -- ++(0,-0.35) -| (pick.north);
\end{tikzpicture}
\caption{한 로터 순서의 탐색 구조. 스크램블러 표는 가운데 로터가 바뀔 때만 다시 만든다. 가장 안쪽
반복에서는 표를 고르고 전류를 퍼뜨리기만 한다.}
\label{fig:tables}
\end{figure}

이 구조는 실제 기계의 움직임과 같은 생각이다. 봄베에서도 한 드럼이 한 바퀴 도는 동안 나머지 두 드럼은
서 있다. 그리고 같은 생각이 FPGA 엔진 설계의 출발점이 된다(\ref{sec:fpga-overview}절). 4로터에서는
그리스 로터와 얇은 반사판이 안쪽 치환에 한 겹 더 들어갈 뿐이고 나머지는 같다.

\subsection{전류의 확산}

\code{propagate}는 전압을 건 선에서 출발해 닿을 수 있는 선을 모두 찾는다. 방법은 작업 목록이다.
새로 통전된 선을 목록에 넣고, 목록에서 하나를 꺼내 그 선과 이어진 선들(메뉴의 선을 따라 이웃 케이블로,
대각선 보드를 따라 대칭 짝으로)을 살핀다. 아직 통전되지 않은 선이면 통전시키고 목록에 넣는다. 목록이
비면 끝난다.

이 계산이 실제 기계의 전류와 같은 답을 내는 이유는, 구하는 것이 ``전압을 건 선과 이어진 선 전체''라는
하나의 집합이기 때문이다. 선을 어떤 순서로 살피든 이 집합은 변하지 않는다. 한 번 통전된 선은 꺼지지
않고, 선은 676개뿐이므로 과정은 반드시 끝난다. 이 성질 덕분에 한 번에 한 가닥씩 살피는 소프트웨어와
모든 선을 동시에 살피는 회로(\ref{ch:rtl}장)가 같은 결과에 도달한다.

구동 중에는 테스트 레지스터의 26선이 모두 통전되는 순간 계산을 그만둔다. 그 위치는 정지가 아닌 것이
이미 확정되었기 때문이다. 대부분의 위치가 여기에 해당하므로 탐색 시간의 대부분이 이 조기 종료로 절약된다.
정지한 위치의 완전한 상태가 필요할 때(그림, 회로 검증 벡터, 체킹)는 조기 종료 없이 끝까지 계산하는
\code{trace}를 쓴다.

\subsection{넘어감 가설}\label{sec:sw-turnover}

모델은 \ref{sec:turnover}절의 문제에 대해 ``크립 위치 $t$에서 가운데 로터가 넘어갔다''는 가설을
지원한다. 가설이 있으면 스크램블러 표를 가운데 위치 $m$과 $m+1$에 대해 두 벌 만들고, 크립 위치가 $t$
이상인 선은 두 번째 벌에서 고른다. 전체 탐색은 $t$의 모든 값에 대해 구동을 반복하므로 일의 양이 메뉴
구간 길이만큼 늘어난다.

예제 전문을 넘어감이 크립 한가운데(크립 위치 12)에 오는 키로 암호화하면, 맞는 로터 순서를 돌려도 가설
없이는 정지가 하나도 나오지 않는다. 가설을 켜면 351,520 위치에서 정지 34회가 나오고 그중 하나가
체킹을 통과하며, 그 정지의 넘어감 가설은 실제 값 12와 일치한다. 가운데 로터의 이중 스테핑이 크립 안에
있는 경우는 다루지 않는다.

\section{체킹 단계}\label{sec:sw-check}

체킹 모듈(\code{check})은 정지 하나를 받아 키 전체와 평문을 내놓거나, 거짓 정지로 기각한다.
그림~\ref{fig:check}의 네 단계로 이루어진다. 앞의 두 단계는 실제 체킹 머신이 하던 일이고, 뒤의 두 단계는
Hut~6에서 손으로 하던 일을 자동화한 것이다.

\begin{figure}[ht]
\centering
\begin{tikzpicture}[x=1cm,y=1cm]
  \node[blkr,minimum width=22mm,minimum height=14mm] (s0) at (0,0) {정지\\{\footnotesize 로터 순서, 위치,}\\{\footnotesize 레지스터}};
  \node[blkb,minimum width=27mm,minimum height=14mm] (s1) at (3.3,0) {1. 스테커 가설\\{\footnotesize 26가지에 전류}};
  \node[blkb,minimum width=27mm,minimum height=14mm] (s2) at (6.7,0) {2. 크립 전체 대조\\{\footnotesize 메뉴 밖 글자 포함}};
  \node[blkg,minimum width=27mm,minimum height=14mm] (s3) at (10.1,0) {3. 링 설정\\{\footnotesize $26^2$가지 시험}};
  \node[blkg,minimum width=27mm,minimum height=14mm] (s4) at (13.5,0) {4. 남은 스테커\\{\footnotesize 언어 통계}};
  \draw[->] (s0) -- (s1); \draw[->] (s1) -- (s2); \draw[->] (s2) -- (s3); \draw[->] (s3) -- (s4);
  \node[note,below=2mm of s1] {케이블에 선이 둘\\살아나면 기각};
  \node[note,below=2mm of s2] {모순이면 기각};
  \node[note,below=2mm of s3] {넘어감이 가정과\\어긋나면 제외};
  \node[note,below=2mm of s4] {키와 평문};
  \draw[decorate,decoration={brace,amplitude=5pt}] (1.9,1.0) -- (8.1,1.0) node[midway,above=6pt,note] {체킹 머신};
  \draw[decorate,decoration={brace,amplitude=5pt}] (8.7,1.0) -- (14.9,1.0) node[midway,above=6pt,note] {Hut 6의 수작업};
\end{tikzpicture}
\caption{체킹 단계의 흐름.}
\label{fig:check}
\end{figure}

\paragraph{1. 스테커 가설.} 테스트 글자의 스테커 짝 후보 26가지 각각에 대해 그 선 하나에 전압을 걸고
전류를 끝까지 퍼뜨린다. 참인 가설은 참인 귀결만 낳으므로 어떤 케이블에도 선이 둘 이상 살아날 수 없다
(\ref{sec:spread}절). 한 케이블에 선이 둘 살아나면 그 가설은 기각된다. 살아남은 가설에서는
케이블마다 살아 있는 선을 읽어 메뉴 글자의 스테커를 얻는다. 봄베가 알려 준 후보만이 아니라 26가지를 모두
시험하는 것은 계산이 싸기 때문이다. 살아남을 수 있는 것은 봄베가 스테커 후보로 알려 준 선뿐이다.

\paragraph{2. 크립 전체 대조.} 메뉴에는 크립의 일부 관계만 들어 있다(스크램블러 수의 제한, 떨어진
덩어리). 이 단계는 크립의 \emph{모든} 글자에 대해 $\Stk(c) = \Scr_k(\Stk(p))$를 확인한다. 한쪽의 스테커를
알면 다른 쪽이 정해지므로 스테커가 더 채워지고, 이미 정해진 값과 어긋나면 기각한다. 더 채워질 것이 없을
때까지 되풀이한다. 거짓 정지의 대부분이 여기서 탈락한다. 예제에서는 157개의 정지 가운데 하나만 통과한다.

\paragraph{3. 링 설정.} 봄베는 코어 위치만 알려 주므로 오른쪽과 가운데 로터의 링 설정은 아직 모른다.
링 설정은 넘어감 시점을 정한다. 모델은 $26\times26$가지 조합을 시험한다. 각 조합에 대해 크립 구간에서의
로터 움직임이 봄베의 가정(넘어감 없음 또는 넘어감 가설)과 맞는지 확인하고, 맞으면 크립 위치에서 전문의
처음으로 로터를 되감아 전문 전체를 복호화한 뒤, 복호문이 독일어다운 정도를 점수로 매긴다. 점수가 가장 높은
조합을 택한다.

여기서 원리적으로 정할 수 없는 것이 있다. 왼쪽 로터의 링은 어떤 넘어감에도 관여하지 않으므로 어떤
값이어도 복호화가 같다. 가운데 로터의 링은 전문 안에서 왼쪽 로터가 움직인 경우에만 정해진다. 모델은
정해지지 않는 링을 \texttt{A}로 두고 시작 위치를 그만큼 조정해, 실제 키와 \emph{등가인} 키를 보고한다.
예제에서 실제 키는 링 \texttt{GMC}, 시작 위치 \texttt{KDC}이고 복원된 키는 링 \texttt{AAC}, 시작 위치
\texttt{ERC}다. 글자는 다르지만 코어 위치(창 글자 $-$ 링)는 둘 다 \texttt{ERA}이고 복호문은 같다.

\paragraph{4. 남은 스테커.} 크립 구간의 평문에도 암호문에도 나오지 않은 글자는 스테커를 논리적으로
정할 수 없다. 이런 글자들에 대해서는 두 글자를 짝지었을 때 전문 전체의 점수가 분명히 오르는 쌍이 있으면
추가하고, 없으면 스테커가 없는 것으로 둔다. 이 부분은 추론이 아니라 추측이므로 결과에 따로
표시한다(\code{unresolved}, \code{guessed}).

\begin{lesson}{복호문 점수는 글자의 빈도가 아니라 글자의 이웃을 봐야 한다}
3단계와 4단계는 복호문이 ``독일어다운 정도''에 기댄다. 처음에는 글자 하나하나의 빈도(독일어에서
\texttt{E}와 \texttt{N}이 흔하다는 사실)로 점수를 매겼다. 이 점수로는 크립에 없는 글자의 스테커를
복원하지 못했다. 글자의 빈도는 글자의 \emph{순서}를 보지 않으므로, 스테커 한 쌍을 잘못 두어 몇 글자가
다른 글자로 바뀌어도 점수가 거의 움직이지 않기 때문이다.

점수를 독일어 \term{바이그램}{연속한 두 글자의 빈도} 통계로 바꾸자 크립에 없는 글자의 스테커까지
복원되었다. 틀린 글자 하나는 앞뒤 두 쌍을 독일어에 없는 조합으로 만들므로 점수에 분명히 드러난다.

그래도 전문의 첫 세 글자가 틀리는 경우가 남았다(\texttt{ANX}가 \texttt{SLE} 또는 \texttt{AGI}로 나왔다).
가운데 로터의 링 설정 후보들이 전문 앞부분에서만 서로 다른 복호문을 주는데, 올바른 복호문에는 군용 전문이
낱말 구분자로 쓰는 \texttt{X}가 들어 있었다. 바이그램 표는 문학 작품에서 만든 것이어서 \texttt{NX},
\texttt{XO} 같은 조합을 거의 본 적이 없었고, 그래서 올바른 답에 낮은 점수를 주었다. 같은 말뭉치에서
낱말 경계를 \texttt{X}로 바꾼 통계를 더해 넣자 이 경우까지 정확해졌다.

교훈은 평범하다. 점수 함수는 실제로 풀려는 문제의 데이터와 같은 성질을 가져야 한다.
\end{lesson}

점수는 연속한 두 글자의 조건부 확률의 로그를 평균한 값이다. 무작위 글자열은 약 $-3.5$, 보통의 독일어
문장은 약 $-2.4$가 나온다. 바이그램 표는 공개 저작물인 독일어 문학 세 편에서 만들었고 생성 스크립트와
함께 \code{data/language/}에 있다.

\section{전체 탐색과 병렬화}

\code{search}는 하루치 키를 찾는 일 전체, 곧 여러 로터 순서에 대해 봄베를 돌리고 모든 정지를 체킹하는
일을 한다. 블레츨리 파크에서 봄베 여러 대가 서로 다른 로터 순서를 나눠 맡던 것에 해당한다.

일은 (로터 순서, 넘어감 가설, 그리스 로터 위치)의 조합으로 쪼개어 CPU의 모든 코어에 나눠 준다.
로터 순서끼리는 서로 영향을 주지 않으므로 그대로 나눌 수 있다. 4로터에서 그리스 로터 위치로 한 번 더
쪼개는 이유는, 로터 순서 하나가 $26^4$ 위치로 커서 로터 순서 몇 개만 돌릴 때에도 모든 코어가 일하게
하기 위해서다. 결과는 복호문 점수가 높은 순으로 정렬해 돌려준다.

\section{Python 조작 인터페이스}

Python 인터페이스는 실제 조작 순서(메뉴 작성, 배선, 구동, 정지 확인, 체킹)를 그대로 따르게 만들었다.
코드~\ref{lst:workflow}는 예제 전문을 푸는 전 과정이다.

\begin{lstlisting}[language=Python,caption={크립에서 평문까지. \code{CIPHER}는 121글자 암호문이다.},label={lst:workflow}]
from christopher import Bombe, Menu, possible_offsets

CRIB = "WETTERVORHERSAGE"
offsets = possible_offsets(CIPHER, CRIB)        # 60 of 106 offsets survive

menu  = Menu.from_crib(CIPHER, CRIB, 28).select(14)   # 14 edges, 14 letters, 1 loop
bombe = Bombe(model="enigma-i", diagonal_board=True).plug_up(menu)

stops = bombe.run(wheel_orders="all")           # 120 wheel orders, 157 stops
for stop in stops:
    for solution in stop.check():               # exactly one stop passes
        print(stop)
        print(solution.wheel_order, solution.rings, solution.positions)
        print(solution.plugboard)
        print(solution.plaintext)
\end{lstlisting}

통과한 정지 하나의 출력은 다음과 같다.

\begin{lstlisting}[numbers=none]
STOP  B III-V-II  drums ETD  test register E: live ABCDEFGHIJKLMNOPQRSTUWXYZ (25/26)
B III-V-II  AAC  ERC
AQ BZ CR DM EV FW GK HN IU JP
ANXOBERKOMMANDODERWEHRMACHTXWETTERVORHERSAGEFUERDIENORDSEEXHEUTEKLARXWIND...
\end{lstlisting}

메뉴 작성부터 체킹까지를 한 번에 하는 \code{search} 함수도 있다. 한국어 튜토리얼
노트북(\code{notebooks/01\_crib\_to\_plaintext.ipynb})이 이 과정을 단계별로 설명하며, 대각선 보드의
유무 비교, 676선 상태의 그림, 넘어감 가설, M4와 실제 전문의 해독을 포함한다.

\section{회로 검증용 벡터}

코어는 회로 검증에 쓸 벡터를 파일로 내보낸다(\code{christopher.vectors.export\_vectors}). 벡터 하나는
메뉴, 로터 순서, 드럼 위치와 그 위치에서의 기대 결과(676선 상태 전체, 테스트 레지스터, 정지 여부)로
이루어진다. 회로는 같은 입력에 대해 같은 676비트를 내야 한다. 벡터의 쓰임은 \ref{sec:rtl-verify}절에서
다룬다.
